\section{Introduction}
Spaced-repetition schedulers decide when a learner should see an item again.
Each decision depends on a predicted probability of recall, so the quality of that predictor drives how much practice is wasted and how much is forgotten.
Recent work compares recall models with each other \cite{settles2016trainable, ye2022stochastic, srsbenchmark}.
A constant recall rate, fitted to training data only, is a floor that uses no memory model at all.
A floor is cheap, and it shows how much a model's structure adds.

We study this question on two datasets that differ a great deal.
Anki gives individually timed reviews of many cards.
Duolingo gives short practice sessions, many learners, and a high base rate of recall.
The two datasets differ in size and in scoring rule, so we report each one on its own terms.

The first result concerns the constant floor on Duolingo.
A smoothed per-user constant has log loss 0.2723.
FSRS-6 (optimised) has 0.3003, and half-life regression (HLR) has 0.4181.
% FACTS: Duolingo, Constant (per user, smoothed) log loss 0.2723; FSRS-6 (optimised) log loss 0.3003; HLR (Settles & Meeder) log loss 0.4181
Our maximum a posteriori (MAP) models v2 and v3 reach 0.2892 and 0.2830, so they also trail the constant.
% FACTS: Duolingo, Ours v2 log loss 0.2892; Ours v3 log loss 0.2830
The Bayesian model v1 reaches 0.2682, which is second only to GBDT, but it improves on the constant by only 0.0040.
% FACTS: Duolingo, Ours (GPU port) minus Constant (per user, smoothed), log loss: diff -0.0040
The v1 model is the same Bayesian family on both datasets: the original implementation on Anki and its GPU port on Duolingo.

The second result is a gradient-boosted decision tree model (GBDT) with features from the FSRS-6 state and from the learner's own history.
On Anki, GBDT has log loss 0.3584, against 0.3794 for FSRS-6 (optimised).
On Duolingo, it has 0.2588, against 0.3003.
% FACTS: Anki, GBDT log loss 0.3584; FSRS-6 (optimised) log loss 0.3794. Duolingo, GBDT log loss 0.2588; FSRS-6 (optimised) log loss 0.3003
On both datasets the log-loss difference is significant under a user-clustered paired bootstrap.
On Anki the difference is $-$0.0210, with 95\% CI [$-$0.0284, $-$0.0144].
On Duolingo it is $-$0.0415, with 95\% CI [$-$0.0423, $-$0.0406].
% FACTS: Anki, GBDT minus FSRS-6 (optimised), log loss: diff -0.0210, 95% CI [-0.0284, -0.0144]. Duolingo, GBDT minus FSRS-6 (optimised), log loss: diff -0.0415, 95% CI [-0.0423, -0.0406]
GBDT also has the lower Brier score and the higher AUC on both datasets.
% FACTS: Anki, GBDT minus FSRS-6 (optimised), Brier: diff -0.0059, 95% CI [-0.0082, -0.0041]; Anki, AUC: diff 0.0285, 95% CI [0.0185, 0.0402]; Duolingo, Brier: diff -0.0078; Duolingo, AUC: diff 0.1273
The Anki calibration difference is not significant, so we do not claim a calibration gain there.
% FACTS: Anki, GBDT minus FSRS-6 (optimised), ECE: diff -0.0042, 95% CI [-0.0132, 0.0012], CI includes 0
% TODO-SRSB: the srs-benchmark comparison of GBDT against RWKV, GRU and FSRS-6 is reported in Section 5; the GBDT numbers there are pending.

The third result is about the Bayesian model v1.
On Anki, v1 is worse than FSRS-6 (optimised), with a log-loss difference of 0.0427 and 95\% CI [0.0327, 0.0554].
% FACTS: Anki, Ours (dynamic, learned lambda) minus FSRS-6 (optimised), log loss: diff 0.0427, 95% CI [0.0327, 0.0554]
Its decay parameter $\lambda$ has not converged on either dataset: in the main runs the split R-hat is 2.44 on Anki and 2.48 on Duolingo, with effective sample sizes of 3.05 and 3.00.
% FACTS: v1 lambda, Anki: split R-hat 2.44, ESS 3.05. v1 lambda, Duolingo (GPU port): split R-hat 2.48, ESS 3.00
The predictions change little under two checks, but $\lambda$ must not be read as an estimate, so the v1 results on both datasets are provisional.
Section 5 treats the Bayesian model as a trade-off between interpretable, uncertainty-aware models and accurate predictors.

The fourth result is a same-session leakage pitfall in the Duolingo traces.
Duolingo records a whole practice session at one timestamp, so a feature built from earlier reviews can see outcomes from the same session.
On the Duolingo development sample, the frozen GBDT's validation log loss was 0.2473 with the leak and 0.2578 after the fix.
% FACTS: Duolingo validation, frozen GBDT, log loss with the same-session leak: 0.2473. Duolingo validation, frozen GBDT, log loss after the fix: 0.2578
We give the grouping rule that removes the leak.

In summary, we contribute:
\begin{itemize}
\item a comparison of recall predictors against constant floors, on two datasets, with one leakage-audited protocol and paired, user-clustered intervals;
\item evidence that a per-user constant beats FSRS-6 and HLR on Duolingo, and that GBDT beats FSRS-6 on both datasets;
\item the same-session leakage pitfall in Duolingo and how to avoid it;
\item a comparison of Bayesian (Markov chain Monte Carlo, MCMC, and MAP) and tree-based predictors, with their trade-offs.
\end{itemize}
Code will be released with the de-anonymised version.
% TODO: release link after review (anonymity).

Section 2 reviews related work.
Section 3 defines the data and the protocol.
Section 4 describes the models.
Section 5 reports the results, Section 6 discusses them, and Section 7 lists the limitations and disclosures.
